\section{Introduction}

Two neural networks can compute the exact same function yet occupy nearly orthogonal positions in weight space --- separated by a mean cosine distance of 1.07, nearly as large as two random unrelated vectors. This is not a corner case: every oracle-constructed sign-flip orbit pair we measure ($N=2{,}500$ total, comprising 5 conditions $\times$ 500 pairs each) exceeds the geometric significance threshold of 0.05 cosine distance. Scaling orbits add a second dimension of the same problem --- mean cosine distance 0.32, again 100\% above threshold.

Weight space learning --- predicting model properties, enabling model merging, supporting continual learning --- depends on encoders that extract meaningful geometric structure from neural network weights. But if two functionally identical networks occupy diametrically opposite positions in weight space, any method operating on raw weights must either learn to ignore this variation or waste representational capacity tracking it. The dominant paradigm, permutation-equivariant encoders \citep{Zhou2023, Navon2023}, addresses the relabeling symmetry of neurons but leaves scaling and sign-flip symmetries --- which theory shows are equally well-defined for ReLU MLPs \citep{Navon2023EquivariantArchitectures} --- empirically uncharacterized in real model zoos.

This paper provides the first empirical quantification of scaling and sign-flip orbit diameters as cosine distances in a real model zoo --- distinct from prior empirical work \citep{Entezari2022, Ainsworth2022} which characterizes permutation symmetry but not scaling or sign-flip orbit geometry.\footnote{Entezari et al.\ \citep{Entezari2022PermutationSymmetry} and Ainsworth et al.\ \citep{Ainsworth2022GitReBasin} characterize the geometry of permutation symmetry orbits in the context of model merging. Schürholt et al.\ \citep{Schurholt2022ModelZoos} characterize statistical properties of model zoo populations. Neither measures the cosine distance diameter of scaling or sign-flip orbits in a real zoo, which is our specific contribution.} We measure the geometric diameter of scaling and sign-flip orbits in the Schürholt MNIST model zoo \citep{Schurholt2022ModelZoos}, a benchmark of 2-layer MLPs with ground-truth property labels. We then probe whether a state-of-the-art equivariant encoder, the Neural Functional Transformer (NFT) \citep{Zhou2023NeuralFunctionalNetworks}, is naturally invariant to these orbits --- and whether explicit canonicalization helps. The answers reveal a more nuanced picture than the theoretical argument alone suggests.

\textbf{NFT is non-invariant to scaling orbits but approximately invariant to sign-flip orbits by construction.} For scaling-orbit pairs, NFT embeddings show within-orbit cosine similarity 0.971, significantly lower than cross-orbit same-property similarity 0.995 (gap$=+0.024$, 95\% CI$=[0.0232, 0.0245]$). For sign-flip pairs, the gap is $-0.0007$ --- 34$\times$ smaller, with within-orbit similarity actually marginally \textit{higher} than cross-orbit. NFT's row-level tokenization operates on magnitude statistics, making it inherently approximately sign-agnostic. This asymmetry was not built by design; it is an emergent property of the architecture.

The natural next step --- canonicalize weights before encoding --- runs into a structural barrier. We show that the majority-sign algorithm, the standard approach to sign-flip canonicalization for $M=2$ ReLU MLPs \citep{Navon2023EquivariantArchitectures}, produces a unique canonical form for only 14.4\% of Schürholt zoo models. For even input dimension $d_{\mathrm{in}}=784$, binomial combinatorics predict an 83\% tie rate per model (observed: 85.6\%), rendering the algorithm non-unique for the overwhelming majority. Property prediction experiments at the available zoo scale ($N=500$) are statistically underpowered --- all condition comparisons yield Spearman $\rho$ confidence intervals of width $\approx 0.6$ --- but directional estimates consistently favor canonicalization ($\Delta\rho = +0.05$--$0.08$ across all tasks). We further caution that full canonicalization (Condition D) did not outperform random normalization (Condition E) at $N=500$; we attribute this to the non-unique sign-flip canonicalization described above, but this attribution remains post-hoc pending adequately-powered follow-up.

\textbf{A note on NFT's current performance.} NFT achieves Spearman $\rho \approx 0.11$ on the Schürholt MNIST zoo at $N=500$, while layer statistics \citep{Unterthiner2020PredictingNN} achieves $\rho \approx 0.9$. This substantial gap is primarily a consequence of underpowered training ($N=400$ training samples) and reflects the invariance problem our work analyzes: NFT devotes representational capacity to tracking symmetry-induced weight variation rather than functional properties. Our canonicalization experiments are designed to close this gap at adequate $N$; the poor baseline $\rho$ motivates, rather than undermines, the invariance analysis.

We make four contributions:

\begin{enumerate}
\item \textbf{Orbit diameter measurement.} First empirical quantification of scaling and sign-flip orbit diameters as cosine distances in a real MLP zoo: scaling mean cosine distance 0.3232 (CI$=[0.3226, 0.3238]$), sign-flip 1.075 (CI$=[1.0705, 1.0789]$), both 100\% above the 0.05 threshold across 2,500 oracle-constructed orbit pairs (5 conditions $\times$ 500 pairs per condition).

\item \textbf{NFT invariance probe.} First quantification of NFT's non-invariance to scaling symmetry (gap$=+0.024$, CI$=[0.0232, 0.0245]$) and emergent approximate invariance to sign-flip symmetry (gap$=-0.0007$) --- an asymmetric result not anticipated by the architecture design.

\item \textbf{Geometric concentration.} Scaling canonicalization increases PCA explained variance ratio from 0.055 to 0.086 at $k=20$ principal components, confirming that canonicalization removes geometric redundancy even when downstream property prediction improvement cannot be detected at $N=500$.

\item \textbf{Structural limitation characterization.} Identification and quantitative characterization of the majority-sign canonicalization failure for even-$d_{\mathrm{in}}$ architectures, with a combinatorial root-cause analysis and a proposed fix (odd $d_{\mathrm{in}}$, alternative tie-breaking, or scaling-only canonicalization).
\end{enumerate}

Together, these contributions transform the argument for symmetry canonicalization in weight space learning from a theoretical expectation into a measurement-grounded framework --- one that specifies exactly which symmetries require canonicalization, which standard algorithms fail and why, and what scale of zoo is needed to detect the downstream effect.

The remainder of the paper is organized as follows. Section~\ref{sec:related} reviews the weight space learning literature and positions our contributions. Section~\ref{sec:method} describes our methodology: oracle orbit construction, NFT invariance probing, and canonicalization implementation. Section~\ref{sec:experiments} presents experimental setup. Section~\ref{sec:results} reports results across all four sub-hypotheses. Section~\ref{sec:discussion} discusses implications, limitations, and the path to adequately-powered follow-up experiments. Section~\ref{sec:conclusion} concludes.
